
\documentclass[a4paper,10pt]{report}

%\usepackage[landscape]{geometry}
\usepackage[cm-default]{fontspec}
\setromanfont{FreeSerif}
\setsansfont{FreeSans}
\setmonofont{FreeMono}
\usepackage[utf8x]{inputenc}
\usepackage{fontspec}
\usepackage{xunicode}
\usepackage{xltxtra}
\usepackage{xgreek}
\usepackage{amsmath}
\usepackage{unicode-math}
\usepackage{ulem}
\usepackage{color}
\usepackage{verbatim}
\usepackage{nopageno}
\usepackage{graphicx}
\usepackage{textpos}
\setlength{\TPHorizModule}{1cm}
\setlength{\TPVertModule}{1cm}
%\usepackage[colorgrid,texcoord]{eso-pic}
\usepackage[outline]{contour}
\usepackage{wrapfig}
\usepackage{url}
\usepackage{color}
\usepackage{tikz}
\usepackage{fancybox,fancyhdr}
\usepackage{subfigure}
\usepackage{pstricks}
\usepackage{epsfig}
\usepackage{multicol}
\usepackage{listings}
\usepackage{enumerate}
\usepackage{hyperref}
\hypersetup{
  bookmarks=true,
  bookmarksopen=true,
  pdfborder=false,
  pdfpagemode=UseNone,
  raiselinks=true,
  pdfhighlight={/P},
  colorlinks,
  citecolor=black,
  filecolor=black,
  linkcolor=black,
  urlcolor=black
}


\usepackage{fontspec}
\usepackage{xunicode}
\usepackage{xltxtra}


% Margins
%\setlength{\textwidth}{24cm}

%\setlength{\voffset}{0in}
%\setlength{\textheight}{6.8in}
% Colors
\definecolor{cbrown}{rgb}{0.49,0.24,0.07}
\definecolor{cmpez}{rgb}{0.92,0.88,0.79}
\definecolor{cmpez2}{rgb}{0.62,0.33,0.34}
\definecolor{cwhite}{rgb}{1,1,1}
\definecolor{cblack}{rgb}{0,0,0}
\definecolor{cred}{rgb}{0.9,0.15,0.15}
\definecolor{clblue}{rgb}{0.57,0.85,0.97}
\definecolor{corange}{rgb}{0.86,0.49,0.18}
\definecolor{clorange}{rgb}{0.95,0.84,0.61}
\definecolor{cyellow}{rgb}{1,0.95,0.16}
\definecolor{cgreen}{rgb}{0.69,0.8,0.31}
\definecolor{clbrown}{rgb}{0.78,0.67,0.43}
\definecolor{cmagenta}{rgb}{0.79,0.05,0.54}
\definecolor{cgray}{rgb}{0.76,0.73,0.63}
\definecolor{bluesite}{rgb}{0.05,0.43,0.67}
\definecolor{redsite}{rgb}{0.90,0.15,0.15}
\definecolor{backsite}{rgb}{0.07,0.07,0.07}
% ----------- from Costas -----------------------------

%---------------   Main Fonts    -----------------%   

% \setmainfont[Mapping=tex-text]{Calibri}
% \setmainfont[Mapping=tex-text]{Times New Roman}
%\setmainfont[Mapping=tex-text]{Droid Serif}
%\setmainfont[Mapping=tex-text]{Cambria}
%\setmainfont[Mapping=tex-text]{cm-unicode}
% \setmainfont[Mapping=tex-text]{Gentium}
% \setmainfont[Mapping=tex-text]{GFS Didot}
% \setmainfont[Mapping=tex-text]{Comic Sans MS}
 \setmainfont[Mapping=tex-text]{Ubuntu}
% \setmainfont[Mapping=tex-text]{Myriad Pro}
% \setmainfont[Mapping=tex-text]{CMU Concrete}
%\setmainfont[Mapping=tex-text]{DejaVu Sans}
%\setmainfont[Mapping=tex-text]{KerkisSans}
%\setmainfont[Mapping=tex-text]{KerkisCaligraphic}
% \setmainfont[Mapping=tex-text]{Segoe Print}
% \setmainfont[Mapping=tex-text]{Gabriola}

%------------------------------------------------%

% ------------ Mathematics Fonts ----------------%
\setmathfont{Asana-Math.ttf}
% -----------------------------------------------%
% Misc
\author{Κασωτάκη Ε. - Σμαραγδάκης Κ.}
\title{www.math24.gr}
% ---------------- Formation --------------------%


\setlength\topmargin{-0.7cm}
\addtolength\topmargin{-\headheight}
\addtolength\topmargin{-\headsep}
\setlength\textheight{26cm}
\setlength\oddsidemargin{-0.54cm}
\setlength\evensidemargin{-0.54cm}
%\setlength\marginparwidth{1.5in}
\setlength\textwidth{17cm}

\RequirePackage[avantgarde]{quotchap}
\renewcommand\chapterheadstartvskip{\vspace*{0\baselineskip}}
\RequirePackage[calcwidth]{titlesec}
\titleformat{\section}[hang]{\bfseries}
{\Large\thesection}{12pt}{\Large}[{\titlerule[0.9pt]}]
%--------------------------------------------------%
\usepackage{draftwatermark}
\SetWatermarkText{www.math24.gr}
\SetWatermarkLightness{0.9}
\SetWatermarkScale{0.6}
%--------------------------------------------------%


\usepackage{xcolor}
\usepackage{amsthm}
\usepackage{framed}
\usepackage{parskip}

\colorlet{shadecolor}{bluesite!20}

\newtheorem*{orismos}{Ορισμός}

\newenvironment{notation}
  {\begin{shaded}\begin{theorem}}
  {\end{theorem}\end{shaded}}




% Document begins
\begin{document}
\pagestyle{fancy}
\fancyhead{}
\fancyfoot{}
\renewcommand{\headrulewidth}{0pt}
\renewcommand{\footrulewidth}{0pt}


\fancyhead[LO,LE]{
 %\textblockcolor{backsite}
 %\begin{textblock}{5}(-2,-0.55)
  %\rule{0cm}{1cm}
 %\end{textblock}
 \textblockcolor{bluesite}
 \begin{textblock}{5}(-1.5,-0.55)
  \rule{0cm}{1cm}
 \end{textblock}
 %\textblockcolor{bluesite}
 %\begin{textblock}{14}(5.5,-0.55)
 % \rule{0cm}{1cm}
 %\end{textblock}
 \begin{textblock}{0}(-1,-0.25)
 \color{cwhite} \begin{Large}www.math24.gr\end{Large}
 \end{textblock}
\begin{textblock}{0}(16,-0.4)
 \color{cwhite} \includegraphics[height=1.5cm]{math24_logo.png}
 \end{textblock}
\textblockcolor{white}
\begin{textblock}{17}(-1,28)
 \color{backsite} \begin{small}Copyright \textcopyright 2012, Κασωτάκη Ε.(ikasotaki@gmail.com) - Σμαραγδάκης Κ.(kesmarag@gmail.com)\end{small}
 \end{textblock}
}
 
\begin{shaded}
\begin{center}
\huge \textbf{Φυλλάδιο Ασκήσεων}\\
%Πρόσθεση ρητών αριθμών
\end{center} 
\textbf{Μαθηματικά Α' Γυμνασίου} \hfill \textbf{Ημερομηνία Παράδοσης : \hspace{2em} }
\subsection*{Ονοματεπώνυμο :\hfill  \hspace{5em}}
\end{shaded}
\vspace{2em}
\begin{itemize}
 \item Πολλαπλάσια ενός αριθμού 
 \item $ΕΚΠ$ δύο ή περισσότερων αριθμών 
 \item Διαιρέτες ενός αριθμού 
 \item $ΜΚΔ$ δύο ή περισσότερων αριθμών 
 \item Πρώτοι αριθμοί 
\end{itemize}

\section*{Θεωρία - ΕΚΠ \hfill \small{}}
\begin{itemize}
 \item \textbf{Πολλαπλάσια} ενός φυσικού αριθμού $α$ είναι όλοι οι αριθμοί που προκύπτουν από τον 
       πολλαπλασιασμό του με όλους τους φυσικούς αριθμούς.
       Δηλαδή τα πολλαπλάσια του αριθμού $α$ είναι: $0,α,2α,3α,4α,5α,\cdots$
 \item \textbf{ΕΚΠ} δύο ή περισσότερων αριθμών (που δεν είναι μηδέν)είναι το μικρότερο από τα κοινά τους πολλαπλάσια.
\end{itemize}

\section*{Θεωρία \hfill \small{}}
\begin{itemize}
 \item Κάθε φυσικός αριθμός διαιρεί τα πολλαπλάσιά του.
 \item Κάθε φυσικός αριθμός που διαιρείται από έναν άλλο είναι πολλαπλάσιό του.
 \item Αν ένας φυσικός διαιρεί έναν άλλο τότε θα διαιρεί και τα πολλαπλάσιά του.
\end{itemize}



\section*{Άσκηση 1  \hfill \small{25 μονάδες}}
Να βρείτε το $ΕΚΠ$ των παρακάτω αριθμών:
\begin{enumerate}[i)]
 \item $5$ και $3$
 \item $7$ και $11$
 \item $2$ και $8$
 \item $3$ και $13$
 \item $5$ και $4$
\end{enumerate}


\section*{Θεωρία - ΜΚΔ \hfill \small{}}
\begin{itemize}
 \item \textbf{Διαιρέτες } ενός φυσικού αριθμού λέγονται όλοι οι αριθμοί που τον διαιρούν.
 \item \textbf{ΜΚΔ} δύο ή περισσότερων αριθμών ονομάζεται ο μεγαλύτερος από τους κοινούς τους διαιρέτες.
\end{itemize}


\section*{Άσκηση 2  \hfill \small{25 μονάδες}}
Να βρείτε το $ΜΚΔ$ των παρακάτω αριθμών:
\begin{enumerate}[i)]
 \item $14$ και $28$
 \item $32$ και $16$
 \item $5$ και $20$
 \item $4$ και $11$
 \item $6$ και $18$
\end{enumerate}

\section*{Θεωρία - Πρώτοι αριθμοί \hfill \small{}}
\begin{itemize}
 \item \textbf{Πρώτοι αριθμοί} είναι οι αριθμοί που έχουν διαιρέτες μόνο τον εαυτό τους και το $1$.\\
       \textbf{Π.χ} ο αριθμός $3$ είναι πρώτος γιατί οι διαιρέτες του είναι οι αριθμοί $1,3$\\
       \textbf{Π.χ} ο αριθμός $6$ δεν είναι πρώτος γιατί οι διαιρέτες του είναι οι αριθμοί $1,2,3,6$
 \item \textbf{Πρώτοι μεταξύ τους} λέγονται δύο αριθμοί όταν ο $ΜΚΔ$ τους είναι το $1$.\\
        Δηλαδή $α,β$ είναι πρώτοι μεταξύ τους αν $ΜΚΔ(α,β)=1$\\
       \textbf{Π.χ} οι αριθμοί $3,10$ είναι πρώτοι μεταξύ τους γιατί  $ΜΚΔ(3,10)=1$\\
       \textbf{Π.χ} οι αριθμοί $3,6$ δεν είναι πρώτοι μεταξύ τους γιατί  $ΜΚΔ(3,6)=3$
\end{itemize}


\section*{Άσκηση 3  \hfill \small{}25  μονάδες}
Να εξετάσετε αν οι παρακάτω αριθμοί είναι πρώτοι:
\begin{enumerate}[i)]
 \item $6$
 \item $7$
 \item $9$
 \item $11$
 \item $15$
\end{enumerate}





\section*{Άσκηση 4  \hfill \small{} 25 μονάδες}
Να εξετάσετε αν οι παρακάτω αριθμοί είναι πρώτοι μεταξύ τους:
\begin{enumerate}[i)]
 \item $3$ και $4$
 \item $4$ και $6$
 \item $6$ και $12$
 \item $10$ και $21$
 \item $13$ και $25$
\end{enumerate}




%\let\thefootnote\relax\footnotetext{Επιμέλεια : Κασωτάκη Ειρήνη (ikasotaki@gmail.com),    Σμαραγδάκης Κώστας (kesmarag@gmail.com)} 


\end{document}
